\chapter{IMPLEMENTASI}

% TODO: Bab implementasi memuat langkah-langkah yang diambil dalam proses
% implementasi berdasarkan rancangan yang telah dibuat, misalnya penulisan
% kode dan integrasi modul. Modul yang diadopsi atau dimodifikasi harus diacu
% sumbernya. Sebutkan alat dan teknologi yang digunakan, seperti bahasa
% pemrograman, kerangka kerja, atau perangkat pengembangan perangkat lunak.

\section{Lingkungan Implementasi}

% TODO: Jelaskan spesifikasi perangkat keras, perangkat lunak, bahasa
% pemrograman, dan kerangka kerja yang digunakan.

Implementasi dilakukan dengan spesifikasi sebagai berikut:

\begin{itemize}
    \item \textbf{Sistem Operasi:} Lorem ipsum dolor sit amet.
    \item \textbf{Perangkat Lunak:} Sed do eiusmod tempor incididunt.
    \item \textbf{Perangkat Keras:} Ut enim ad minim veniam.
    \item \textbf{Bahasa Pemrograman:} Lorem ipsum dolor sit amet.
    \item \textbf{Kerangka Kerja:} Sed do eiusmod tempor incididunt.
\end{itemize}

\section{Implementasi Metode/Model}

% TODO: Jelaskan implementasi metode/model sesuai rancangan. Sertakan kode
% yang relevan beserta penjelasan bagian-bagian pentingnya. Untuk modul yang
% diadopsi atau dimodifikasi, cantumkan sumbernya.

\subsection{Modul A}

% TODO: Jelaskan implementasi modul A.

Lorem ipsum dolor sit amet, consectetur adipiscing elit. Sed do eiusmod tempor incididunt ut labore et dolore magna aliqua.

\begin{lstlisting}[language=python, caption={Contoh Implementasi Modul A}]
# TODO: Ganti dengan kode implementasi Anda
def example_function():
    print("Hello, World!")
    return True
\end{lstlisting}

\subsection{Modul B}

% TODO: Jelaskan implementasi modul B.

Lorem ipsum dolor sit amet, consectetur adipiscing elit. Sed do eiusmod tempor incididunt ut labore et dolore magna aliqua.

\section{Integrasi Modul}

% TODO: Jelaskan bagaimana modul-modul diintegrasikan menjadi satu kesatuan
% sistem yang berfungsi.

Lorem ipsum dolor sit amet, consectetur adipiscing elit. Sed do eiusmod tempor incididunt ut labore et dolore magna aliqua. Ut enim ad minim veniam, quis nostrud exercitation ullamco laboris nisi ut aliquip ex ea commodo consequat.

\section{Pengujian Hasil Implementasi dan Debugging}

% TODO: Deskripsikan pendekatan pengujian hasil implementasi sesuai teknik
% yang disampaikan pada rancangan, paparkan hasil tes yang diperoleh, serta
% jelaskan proses pelacakan kesalahan dan penyelesaiannya selama pengujian.

Lorem ipsum dolor sit amet, consectetur adipiscing elit. Sed do eiusmod tempor incididunt ut labore et dolore magna aliqua.

% TODO: Lampirkan dokumentasi kode lengkap dengan komentar pada bagian
% lampiran agar mudah dipahami dan dipelihara.
